\documentclass[10pt]{article}

% ---------- Document Identity (update these only) ----------
\newcommand{\DocStage}{}
\newcommand{\DocTitle}{Your Road to Tier One SOF}
\newcommand{\ProgramName}{182-Week Civilian-to-Selection Readiness Pipeline}
\newcommand{\ProgramSubtitle}{}

% ---------- Page ----------
\usepackage[legalpaper,left=0.5in,right=0.5in,top=0.6in,bottom=0.45in]{geometry}

% ---------- Typography ----------
\usepackage[T1]{fontenc}
\usepackage[utf8]{inputenc}
\usepackage{libertinus}
\usepackage{microtype}
\usepackage{parskip}
\usepackage{ragged2e}
\usepackage{textcomp}
\AtBeginDocument{\color{black}}
\AtBeginDocument{\pagecolor{white}}
\AtBeginDocument{\justifying}

% ---------- Landscape pages ----------
\usepackage{pdflscape}

% ---------- Color ----------
\usepackage[table]{xcolor}
\definecolor{StageBlue}{HTML}{1F4E79}
\definecolor{StageBlueDark}{HTML}{163A5A}
\definecolor{LightGray}{HTML}{F2F4F7}
\definecolor{MidGray}{HTML}{D9DEE5}
\definecolor{RuleGray}{HTML}{C7CED8}
\definecolor{HeaderInk}{HTML}{000000}
\definecolor{Ink}{HTML}{000000}

% ---------- Utility ----------
\usepackage{enumitem}
\sloppy
\setlength{\emergencystretch}{2em}

% ---------- Boxes / UI ----------
\usepackage[most]{tcolorbox}

\tcbset{
  charterTitle/.style={
    enhanced,
    colback=StageBlue!4,
    colframe=StageBlue,
    boxrule=1.25pt,
    arc=3pt,
    left=16pt,right=16pt,top=14pt,bottom=14pt,
    borderline north={3.2pt}{0pt}{StageBlue},
    borderline west={4pt}{0pt}{StageBlue},
    borderline south={1.25pt}{0pt}{StageBlue},
    drop shadow={black!25!white},
  },
  charterRule/.style={
    enhanced,
    colback=LightGray,
    colframe=RuleGray,
    boxrule=0.85pt,
    arc=3pt,
    left=12pt,right=12pt,top=10pt,bottom=10pt,
    borderline west={3pt}{0pt}{StageBlue},
  },
  charterBadge/.style={
    enhanced,
    colback=StageBlue,
    coltext=white,
    colframe=StageBlueDark,
    boxrule=0.6pt,
    arc=2pt,
    left=6pt,right=6pt,top=3pt,bottom=3pt,
  }
}
\newtcbox{\Badge}{nobeforeafter,charterBadge}

% ---------- Lists ----------
\setlist[itemize]{leftmargin=1.5em, itemsep=0.25em, topsep=0.25em}
\setlist[enumerate]{leftmargin=1.8em, itemsep=0.25em, topsep=0.25em}

% ---------- TOC ----------
\usepackage{tocloft}
\setcounter{tocdepth}{3}
\setlength{\cftbeforesecskip}{3pt}
\setlength{\cftbeforesubsecskip}{2pt}
\setlength{\cftbeforesubsubsecskip}{0pt}
\renewcommand{\cftsecfont}{\bfseries}
\renewcommand{\cftsecpagefont}{\bfseries}
\renewcommand{\cftsubsecfont}{\normalfont}
\renewcommand{\cftsubsecpagefont}{\normalfont}
\renewcommand{\cftsubsubsecfont}{\normalfont}
\renewcommand{\cftsubsubsecpagefont}{\normalfont}
\setlength{\cftsecindent}{0pt}
\setlength{\cftsubsecindent}{1.25em}
\setlength{\cftsubsubsecindent}{2.8em}
\setlength{\cftsecnumwidth}{2.1em}
\setlength{\cftsubsecnumwidth}{2.8em}
\setlength{\cftsubsubsecnumwidth}{3.6em}

% Remove the built-in TOC heading so only the custom heading is shown
\makeatletter
\renewcommand{\tableofcontents}{\@starttoc{toc}}
\makeatother

% ---------- Header/Footer: page number only ----------
\usepackage{fancyhdr}
\pagestyle{fancy}
\fancyhf{}
\renewcommand{\headrulewidth}{0pt}
\renewcommand{\footrulewidth}{0pt}
\fancyfoot[C]{\thepage}

% ---------- Section formatting ----------
\usepackage{titlesec}

\titleformat{\section}
  {\Large\bfseries\color{Ink}}
  {\thesection}{0.6em}{}
  [\vspace{0.25em}\textcolor{StageBlue}{\rule{\linewidth}{0.8pt}}]

\titleformat{\subsection}
  {\large\bfseries\color{Ink}}
  {\thesubsection}{0.6em}{}

\titleformat{\subsubsection}
  {\normalsize\bfseries\color{Ink}}
  {\thesubsubsection}{0.6em}{}

\titlespacing*{\section}{0pt}{0.95em}{0.35em}
\titlespacing*{\subsection}{0pt}{0.65em}{0.25em}
\titlespacing*{\subsubsection}{0pt}{0.55em}{0.2em}

% ---------- Table engine ----------
\usepackage{tabularray}
\UseTblrLibrary{booktabs}

% ---------- tabularray: GLOBAL table treatment ----------
\SetTblrInner{
  cells = {fg=black, bg=white, valign=t},
}

% ---------- Header helpers ----------
\newcommand{\Hdr}[1]{\shortstack[c]{\strut #1\strut}}

% ---------- Lines ----------
\newcommand{\TitleLine}{\textcolor{StageBlue}{\rule{0.98\linewidth}{0.95pt}}}
\newcommand{\SubTitleLine}{\textcolor{RuleGray}{\rule{0.98\linewidth}{0.65pt}}}

% ---------- Continued on next page ----------
\newcommand{\ContinuedOnNextPageInline}{%
  \par
  \vfill
  \noindent\textcolor{RuleGray}{\rule{\linewidth}{1pt}}\vspace{-2pt}
  \noindent\hfill\textit{Continued on next page}\par\vspace{-10pt}
  \noindent\textcolor{RuleGray}{\rule{\linewidth}{1pt}}
  \clearpage
}

% ---------- PDF hyperlinks ----------
\usepackage[hidelinks]{hyperref}
\hypersetup{
  colorlinks=false,
  pdfborder={0 0 0},
  linktoc=all,
  pdftitle={\DocTitle},
  pdfauthor={\ProgramName}
}

% =========================================================
\begin{document}

\subsection*{Phase 02 Card Header}
\phantomsection
\addcontentsline{toc}{subsection}{Phase 02 Card Header}

\begin{itemize}
\item Phase \textbf{ID}: Phase 02
\item Phase Name: Maximal Strength Development
\item Duration weeks: 12
\item Version: v1.0
\item Stage Alignment: Stage 5
\item Governing Status: Draft — Non-Deployable
\end{itemize}

\subsection*{1. Phase Mission}
\phantomsection

Build strength reserves under submax, RPE-capped proxy posture while maintaining \textbf{CAL} battery standardization and continuing \textbf{RUN} development using protected-window evidence only. Phase 02 must increase the strength ceiling without creating durability breakdown, evidence invalidity, or unsafe stacking of major tests.

\subsection*{2. Phase Purpose}
\phantomsection

\begin{itemize}
\item Raise the candidate’s strength reserve so later phases can convert it into power-endurance and work capacity without exceeding tissue tolerance.
\item Transition from Phase 01’s “build tolerance + standardize \textbf{CAL}/\textbf{RUN} evidence” into a phase where external loading is treated as a controlled capability with strict comparability controls.
\item Maintain and refine \textbf{CAL} battery execution standards (artifact-valid) so Phase 03 can increase density without redefining counting rules.
\item Maintain \textbf{RUN} evidence continuity without forcing increased impact risk; Phase 02 uses the 2-mile construct (\textbf{C9} / \textbf{MET-2B-015}) as the primary \textbf{RUN} verification inside protected windows.
\item Establish a repeatable \textbf{STR} proxy evidence chain (\textbf{C17}, only if finalized and deployable, / \textbf{MET-2B-027}-031) that is stable in implement identity, setup, and logging so results are comparable across windows and phases.
\item Preserve adherence and validity infrastructure: Phase 02 only “counts” when evidence is admissible under Stage 1.F and Stage 2, and when protected-window posture is respected under Stage 3.
\end{itemize}

\subsection*{3. Entry Criteria}
\phantomsection

\begin{itemize}
\item Prior phase disposition: Phase 01 end gate outcome is \textbf{ADVANCE}, documented under Stage 1.F gate doctrine using admissible evidence only.
\item Session admissibility readiness (Stage 1.F default decision window posture):
\begin{itemize}
\item The two most recent complete training weeks preceding Phase 02 Week 1 demonstrate at least 5 of 6 admissible sessions per week on average.
\item Outside medical issues, no more than 1 \textbf{MODIFIED} session per week is admissible.
\item Any \textbf{INVALID} session inside the decision window disqualifies \textbf{ADVANCE} posture and forces \textbf{HOLD} until corrected.
\end{itemize}
\item Compliance prerequisites:
\begin{itemize}
\item \textbf{MET-2B-044} (Session Validity Tag) is assigned same day for every session in the decision window.
\item \textbf{MET-2B-045} (Cardio Minimum Standard Met) is Yes for each admissible session in the decision window.
\item \textbf{MET-2B-046} (Weekly Admissible Session Compliance Percent) supports the 5-of-6 admissibility posture.
\end{itemize}
\item Durability prerequisites:
\begin{itemize}
\item \textbf{MET-2B-036} (Gait-Altering Pain Flag) is No across the decision window.
\item \textbf{MET-2B-037} (Foot Hotspot or Blister Flag) does not indicate an active blister state that would predictably compromise \textbf{RUN} mechanics or external loading stability.
\end{itemize}
\item Test readiness posture:
\begin{itemize}
\item \textbf{CAL} battery artifact workflow is executable (camera angle, timer visibility, counting auditable) so \textbf{C4}–\textbf{C7} can be \textbf{VALID} when scheduled.
\item A stable \textbf{RUN} environment exists for \textbf{C9}: either a verified outdoor route identity or a stable treadmill identity with required settings logging.
\item A stable \textbf{STR} proxy environment exists for \textbf{C17} only if \textbf{C17} is finalized and deployable: implement identity, loading method, and setup can be held stable enough to make results comparable.
\end{itemize}
\item Default posture if any prerequisite is not met: \textbf{HOLD} at Phase transition, correct prerequisites, and reslot baseline/protected window use per Stage 3; do not “start anyway” under compromised evidence posture.
\end{itemize}

\subsection*{4. Operating Constraints}
\phantomsection

\begin{itemize}
\item Six sessions per week are mandatory across Phase 02. Risk is managed by dose, modality, and intensity substitution—not by collapsing weekly frequency.
\item Every session \textbf{MUST} include, in this exact order:
\begin{itemize}
\item Safety self-check first
\item Warm-up
\item Mobility
\item Main work
\item Cardio
\item Cooldown
\item Recovery notes and any required post-session notes per Stage 1.F
\end{itemize}
\item Cardio minimum standard: minimum 30 minutes continuous per session. If continuous work cannot be maintained safely, downshift intensity and/or substitute to a lower-risk modality to preserve continuous duration; aborted sessions are tagged per Stage 1.F and are not progression-eligible.
\item 10,000 steps per day is a global requirement and is treated as baseline activity compliance evidence; steps do not replace session structure requirements.
\item Gate evidence posture: gate decisions rely on admissible evidence only. Tests must be \textbf{VALID}; \textbf{INVALID} tests provide no gate credit and require reslot to the next protected window per Stage 3.
\item No underwater or breath-hold content is introduced or implied as routine. Underwater governance constructs exist only as Stage 2 instruments for later phases and are not scheduled here.
\item No ruck time trials are required as gate evidence in Phase 02 by default; any ruck exposure remains training-focused and must not be treated as equivalent evidence for \textbf{RUN} or \textbf{STR} constructs.
\end{itemize}

\subsection*{5. Primary Adaptations and Physiological Focus}
\phantomsection

\begin{itemize}
\item Strength reserve expansion: increase load-handling capacity under submax posture, emphasizing repeatable technique and controlled effort ceilings.
\item Neuromuscular efficiency: improved motor unit recruitment and pattern stability in major movement patterns without failure-driven degradation.
\item Connective tissue tolerance: gradual exposure to higher tension under controlled mechanics, reducing the likelihood of sudden overload injuries.
\item Trunk/bracing and posture robustness: improved ability to maintain stable positions under load, supporting later work capacity and ruck tolerance phases.
\item Maintenance of aerobic base: preserve continuous Cardio throughput to support recovery capacity and prevent detraining while strength emphasis rises.
\item Durability governance: maintain stable gait mechanics and foot integrity under increased total training stress; explicitly prevent drift into pain-compensated patterns.
\end{itemize}

\subsection*{6. Primary Modalities}
\phantomsection

Domain emphasis (primary):

\begin{itemize}
\item \textbf{STR}: controlled external loading and strength pattern development; \textbf{STR} proxy testing via \textbf{C17} in protected windows.
\item \textbf{CAL}: continued standardization and maintenance of the \textbf{CAL} battery; artifact-valid outcomes in protected windows.
\item \textbf{RUN}: continued development with controlled, test-week-only time-trial posture using \textbf{C9} (2-mile) as the primary Phase 02 \textbf{RUN} verification.
\end{itemize}

Domain support (continuous monitoring and constraints):

\begin{itemize}
\item \textbf{DURABILITY}: pain severity and disqualifier flags constrain whether high-cost tests are attempted.
\item \textbf{RECOVERY}: sleep/fatigue/soreness and session RPE constrain scheduling and protect validity.
\item \textbf{BODYCOMP}: trend capture continues as gate context and risk control, not as a performance proxy.
\end{itemize}

De-emphasized or prohibited as gate constructs in this phase:

\begin{itemize}
\item \textbf{RUCK}: not a default gate domain in Phase 02; do not treat ruck exposure as equivalent evidence for \textbf{RUN} or \textbf{STR}.
\item \textbf{SWIM}: not scheduled as a default major test event in Phase 02 unless explicitly inserted under cap/spacing feasibility and verified pool conditions; no underwater constructs.
\item \textbf{AER} time trials: not used as \textbf{RUN} substitutes for gate purposes; may exist later as supporting constructs but are not central to Phase 02 gate posture.
\end{itemize}

\clearpage
\begin{landscape}

\subsection*{7. Weekly Structure}
\phantomsection

\begingroup
\scriptsize
\setlength{\tabcolsep}{2pt}
\renewcommand{\arraystretch}{1.12}

\begin{longtblr}{
width=\linewidth,
rowhead=1,
colspec = {
Q[l,wd=0.60in]
X[l,1.80]
X[l,1.25]
X[l,2.75]
X[l,2.60]
},
hlines,
vlines,
cells = {hsep=6pt, vsep=6pt, halign=l, valign=t},
cell{1-Z}{1-5} = {fg=black},
row{odd} = {bg=white},
row{even} = {bg=LightGray},
row{1} = {bg=StageBlue!15, fg=black, font=\bfseries, halign=l, valign=t},
}
\Hdr{Session #} &
\Hdr{Primary Focus} &
\Hdr{Main Work Domains} &
\Hdr{Cardio Modality/Intent} &
\Hdr{Notes and Risk Controls} \\

1 &
\textbf{STR} emphasis lower-body patterns, submax posture &
\textbf{STR}; \textbf{DURABILITY} &
Modality: step-based locomotion (walk or treadmill walk); Minutes: 30–45; RPE+Talk Test: RPE 3–5, full sentences; Continuity: continuous planned &
Technique-first and bracing-first posture. Any mechanics-altering pain terminates the provoking exposure and triggers same-day regression. Preserve low-impact Cardio; do not chase fatigue. \\

2 &
\textbf{STR} emphasis upper-body patterns + \textbf{CAL} quality maintenance &
\textbf{STR}; \textbf{CAL}; \textbf{DURABILITY} &
Modality: low-impact machine or walk (walk, bike, or elliptical); Minutes: 30–45; RPE+Talk Test: RPE 3–5, full sentences; Continuity: continuous planned &
Shoulder history requires conservative end-range posture; prioritize controlled range and scapular stability. Avoid failure-driven work. \\

3 &
\textbf{RUN} mechanics and aerobic base continuity (no time-trial posture) &
\textbf{RUN}; \textbf{RECOVERY}; \textbf{DURABILITY} &
Modality: walk-jog exposure only if mechanics are stable; otherwise walk/bike/elliptical; Minutes: 30–50; RPE+Talk Test: RPE 3–5, full sentences; Continuity: continuous planned &
This is not a test day. If Grand Forks conditions are unsafe (ice/visibility), substitute to indoor modality; document environment. No intensity chasing. \\

4 &
\textbf{STR} supplemental patterns + trunk/bracing quality &
\textbf{STR}; \textbf{DURABILITY}; \textbf{RECOVERY} &
Modality: step-based locomotion or bike; Minutes: 30–45; RPE+Talk Test: RPE 3–4, full sentences; Continuity: continuous planned &
Keep total stress density controlled. If sleep/fatigue trends are degraded, this session becomes Minimum Viable Session posture while preserving structure and Cardio continuity. \\

5 &
\textbf{CAL} battery technique integrity (non-test), mobility bias &
\textbf{CAL}; \textbf{RECOVERY}; \textbf{DURABILITY} &
Modality: walk/elliptical/bike; Minutes: 30–45; RPE+Talk Test: RPE 3–4, full sentences; Continuity: continuous planned &
\textbf{CAL} work remains strict-standard aware; do not chase maximal attempts outside protected windows. Prepare artifact setup and standards consistency approaching test weeks. \\

6 &
Cardio-primary low-impact endurance exposure + recovery emphasis &
\textbf{RECOVERY}; \textbf{DURABILITY} &
Modality: walk (outdoor if safe) or treadmill, bike, or elliptical; Minutes: 45–60; RPE+Talk Test: RPE 3–5, full sentences to short sentences; Continuity: continuous planned &
This is the primary Cardio-density day. Protect foot integrity: if hotspots emerge, substitute to bike/elliptical and document. No breathless intensity chasing. \\

\end{longtblr}

\endgroup

\subsection*{8. Progression Rules}
\phantomsection

Global progression controls (Phase 02 safe posture):

\begin{itemize}
\item No failure-driven training posture. Technique quality is the gate: terminate sets/efforts before form degradation.
\item One progression lever at a time rule: do not increase more than one primary lever in a single week across (a) \textbf{STR} intensity exposure, (b) overall \textbf{STR} volume density, (c) \textbf{RUN} impact exposure, (d) longest Cardio-primary duration.
\item Protected-window no-novelty posture: in the lead-in to Weeks 4/8/12 and during those windows, avoid tool changes, footwear changes, route/machine changes, and artifact workflow changes unless forced; forced changes are comparability breaks and bias toward conservative interpretation and/or reslot.
\item Regression precedence: regress ROM, leverage, complexity, and intensity before reducing weekly frequency. Weekly frequency remains 6 sessions unless \textbf{STOP} \textbf{NOW} or \textbf{MEDICAL} \textbf{HOLD} posture requires deviation.
\end{itemize}

\textbf{STR} progression rules (non-programming; caps and safeguards):

\begin{itemize}
\item \textbf{C17} posture is mandatory for any \textbf{STR} proxy evidence: submax, implement-agnostic, RPE-capped; no true \textbf{1-RM} attempts.
\item \textbf{STR} progression is permitted only when \textbf{DURABILITY} markers do not signal mechanics drift:
\begin{itemize}
\item \textbf{MET-2B-036} (gait-altering pain flag) must remain No.
\item \textbf{MET-2B-037} (foot hotspot/blister) must not indicate worsening skin integrity that predicts compensations.
\end{itemize}
\item If pain alters mechanics during \textbf{STR} work (limp, compensation, unstable joint position), terminate that exposure immediately and substitute to a lower-risk variant. The session is tagged \textbf{MODIFIED} with \textbf{SESSION-RC} and notes.
\item \textbf{STR} proxy comparability control: implement identity and loading method must remain stable across test events used for tier claims; tool/implement drift forces conservative interpretation until stabilized and replicated.
\end{itemize}

\textbf{CAL} progression rules (non-programming; integrity posture):

\begin{itemize}
\item \textbf{CAL} battery results are gate-admissible only when artifact evidence meets Stage 2.C requirements; otherwise tests are \textbf{INVALID} for gate use and must be reslotted.
\item Outside protected windows, \textbf{CAL} work is technique and capacity development only; do not treat work-week maximal attempts as gate-grade evidence.
\end{itemize}

\textbf{RUN} progression rules (Phase 02 impact governance):

\begin{itemize}
\item \textbf{RUN} time-trial evidence (\textbf{C9} / \textbf{MET-2B-015}) is produced only in protected windows.
\item If Grand Forks conditions compromise safety/validity, treadmill is admissible only with full treadmill identity/settings logging per Stage 2 posture; otherwise reslot.
\item \textbf{RUN} progression is not permitted to “make up” missed tests. Missed or \textbf{INVALID} tests trigger reslot discipline, not compression.
\end{itemize}

Cardio progression rules (global constraints embedded):

\begin{itemize}
\item Cardio is completed every session as a continuous block meeting the minimum standard; intensity is controlled by RPE and talk test.
\item Step-based locomotion is the default when foot integrity and environment safety allow; bike/elliptical substitutions are the default risk-reduction tool when foot or joint risk signals rise.
\end{itemize}

\subsection*{9. Deload and Test Placement}
\phantomsection

\begin{itemize}
\item Phase duration: 12 weeks.
\item Default work-to-deload posture: 3:1 with protected Deload+Test windows at:
\begin{itemize}
\item Week 1 (baseline) — protected window for baseline capture and infrastructure checks.
\item Week 4 (mid) — protected Deload+Test window.
\item Week 8 (mid) — protected Deload+Test window.
\item Week 12 (end) — protected Deload+Test end gate window.
\end{itemize}
\item Major event caps and spacing doctrine apply inside each protected week:
\begin{itemize}
\item Ideal cap: 2 major test events.
\item Hard cap: 3 major test events.
\item Any two major test events must respect 72-hour spacing (start-to-start).
\end{itemize}
\item When cap/spacing cannot be protected inside the week, the correct posture is partial completion and reslot to the next protected window (often Phase 03 Week 1 baseline), not compression into build weeks.
\end{itemize}

\end{landscape}

\clearpage
\begin{landscape}

\subsection*{10. Test Battery}
\phantomsection

\begingroup
\scriptsize
\setlength{\tabcolsep}{2pt}
\renewcommand{\arraystretch}{1.12}

\begin{longtblr}{
width=\linewidth,
rowhead=1,
colspec = {
X[l,1.20]
X[l,3.10]
Q[l,wd=0.85in]
X[l,2.85]
X[l,2.00]
},
hlines,
vlines,
cells = {hsep=6pt, vsep=6pt, halign=l, valign=t},
cell{1-Z}{1-5} = {fg=black},
row{odd} = {bg=white},
row{even} = {bg=LightGray},
row{1} = {bg=StageBlue!15, fg=black, font=\bfseries, halign=l, valign=t},
}
\Hdr{Test Timing} &
\Hdr{Test \textbf{ID}/Name} &
\Hdr{Domain} &
\Hdr{Validity Conditions} &
\Hdr{Pass Threshold Stage 2 Tier} \\

Week 1 (baseline) &
\textbf{CAL} battery: \textbf{C4} Push-Ups (\textbf{MET-2B-010}) + \textbf{C5} Sit-Ups (\textbf{MET-2B-011}) + \textbf{C6} Pull-Ups (\textbf{MET-2B-012}, if viable) + \textbf{C7} Plank (\textbf{MET-2B-013}) &
\textbf{CAL} &
Stage 2 \textbf{CAL} validity posture enforced; compliant artifact evidence required for gate-validity; standardized setup and auditable timing/counting; any missing/non-compliant artifact = \textbf{INVALID} for gate use &
Baseline capture; no phase \textbf{ADVANCE} claim from W1 alone; used to establish standards and support replication posture \\

Week 1 (baseline) &
\textbf{STR} proxy baseline session: \textbf{C17} Strength Proxy Test only if \textbf{C17} is finalized and deployable (\textbf{MET-2B-027}; \textbf{MET-2B-028}; \textbf{MET-2B-029}; \textbf{MET-2B-030}; \textbf{MET-2B-031}) &
\textbf{STR} &
Submax, RPE-capped per \textbf{C17}; implement/setup identity logged; warm-up logged; no stop-rule violations; incomplete \textbf{MEASURE}/identity fields = \textbf{INVALID} for gate use &
Baseline capture; no phase \textbf{ADVANCE} claim from W1 alone; establishes comparability baseline for later proxy events \\

Week 4 (mid) &
\textbf{STR} proxy session: \textbf{C17} (\textbf{MET-2B-027}–031) &
\textbf{STR} &
Same as above; treat as a major test event; enforce 72-hour spacing from \textbf{CAL} battery &
Check-in against \textbf{INTERMEDIATE} trajectory; informational for gate planning unless replicated and stable \\

Week 4 (mid) &
\textbf{CAL} battery: \textbf{C4}–\textbf{C7} (\textbf{MET-2B-010}–013; \textbf{C6} if viable) &
\textbf{CAL} &
Artifact evidence required; counting/timing must be auditable; missing artifact = \textbf{INVALID} for gate use &
Check-in against \textbf{INTERMEDIATE} trajectory; informational for gate planning unless replicated and stable \\

Week 8 (mid) &
\textbf{RUN} \textbf{TT}: \textbf{C9} 2-Mile Run Time Trial (\textbf{MET-2B-015}) &
\textbf{RUN} &
Route identity (\textbf{ROUTE}-2E-###) or treadmill identity (\textbf{TM}-2E-###) logged; treadmill grade/settings logged when used; warm-up logged; no stop-rule violations; unsafe environment or missing \textbf{MEASURE} fields = \textbf{INVALID} for gate use &
Check-in against \textbf{INTERMEDIATE} trajectory; supports replication chain for end gate \\

Week 8 (mid) &
\textbf{CAL} battery: \textbf{C4}–\textbf{C7} (\textbf{MET-2B-010}–013; \textbf{C6} if viable) &
\textbf{CAL} &
Artifact evidence required; enforce spacing from \textbf{RUN} \textbf{TT} within the protected week &
Check-in against \textbf{INTERMEDIATE} trajectory; supports replication chain for end gate \\

Week 12 (end gate) &
\textbf{RUN} \textbf{TT}: \textbf{C9} 2-Mile Run Time Trial (\textbf{MET-2B-015}) &
\textbf{RUN} &
Same \textbf{RUN} validity posture; treat as decisive gate-grade attempt; reslot if unsafe or compromised &
\textbf{INTERMEDIATE} tier required for Phase 03 entry \\

Week 12 (end gate) &
\textbf{CAL} battery: \textbf{C4}–\textbf{C7} (\textbf{MET-2B-010}–013; \textbf{C6} if viable) &
\textbf{CAL} &
Artifact evidence required; treat as decisive gate-grade attempt; reslot if artifact cannot be satisfied &
\textbf{INTERMEDIATE} tier required for Phase 03 entry \\

Week 12 (end gate; staged/cap-limited) &
\textbf{STR} proxy session: \textbf{C17} (\textbf{MET-2B-027}–031) &
\textbf{STR} &
Executed only if hard cap and 72-hour spacing can be satisfied inside Week 12; if spacing cannot be protected, reslot \textbf{STR} proxy to Phase 03 Week 1 baseline window &
\textbf{INTERMEDIATE} tier required if used as Phase 02 exit evidence; if reslotted, Phase 02 exit uses the most recent \textbf{VALID} Phase 02 \textbf{STR} proxy session (Week 4 anchor) plus \textbf{RUN}+\textbf{CAL} gate evidence \\

\end{longtblr}

\endgroup

\end{landscape}
\clearpage

\subsection*{11. Exit Standards}
\phantomsection

Phase 02 end gate (Week 12 Deload+Test) produces the decision for entry into Phase 03. Gate decision rules are constrained by Stage 1.F admissibility doctrine and Stage 3 scheduling doctrine.

A) Evidence admissibility and eligibility (Stage 1.F default rules; enforced here)

\begin{itemize}
\item Decision window: the two most recent complete training weeks preceding the Week 12 gate week.
\item \textbf{ADVANCE} eligibility: at least 5 of 6 admissible sessions per week on average across the decision window.
\item \textbf{MODIFIED} cap: outside medical issues, no more than 1 \textbf{MODIFIED} session per week is admissible.
\item Binary disqualifier: any \textbf{INVALID} session within the decision window disqualifies \textbf{ADVANCE} and forces \textbf{HOLD}.
\item Required tests for the gate must be \textbf{VALID}; any required test tagged \textbf{INVALID} defaults \textbf{HOLD} for that metric until a valid retest exists.
\item Dependencies for Phase 03 entry must be validated; risk flags and trends must be reviewed, including durability disqualifiers.
\end{itemize}

B) Required gate-grade performance standards (Week 12)

\begin{itemize}
\item \textbf{RUN} gate construct:
\begin{itemize}
\item \textbf{C9} / \textbf{MET-2B-015} is executed as a \textbf{VALID} test inside Week 12 protected window.
\item Phase 02 exit requires \textbf{INTERMEDIATE} tier attainment on \textbf{C9} / \textbf{MET-2B-015} under \textbf{VALID} conditions, subject to Stage 2 replication posture.
\end{itemize}
\item \textbf{CAL} gate construct:
\begin{itemize}
\item \textbf{CAL} battery \textbf{C4}–\textbf{C7} (\textbf{C6} if viable) is executed as \textbf{VALID} for gate use, including compliant artifact evidence.
\item Phase 02 exit requires \textbf{INTERMEDIATE} tier attainment for the \textbf{CAL} battery constructs under \textbf{VALID} conditions, subject to Stage 2 replication posture.
\end{itemize}
\item \textbf{STR} phase-identity construct (strength reserve evidence):
\begin{itemize}
\item At least one \textbf{VALID} \textbf{C17} \textbf{STR} proxy session exists within Phase 02 only if \textbf{C17} is finalized and deployable (Week 4 is the anchor by default).
\item Phase 02 exit requires \textbf{INTERMEDIATE} tier attainment for the \textbf{C17} proxy posture when evaluated under the most recent valid evidence chain; if Week 12 \textbf{STR} proxy is reslotted due to cap/spacing, the most recent \textbf{VALID} Phase 02 \textbf{STR} proxy session is used for Phase 02 exit review.
\item No \textbf{STR} evidence is admissible for gate claims if implement identity and setup method is not logged.
\end{itemize}
\end{itemize}

C) Disqualifiers for \textbf{ADVANCE} to Phase 03 (Phase 02-specific application of Stage 1.F posture)

\begin{itemize}
\item Any \textbf{INVALID} session in the decision window.
\item Any required gate test (\textbf{RUN} \textbf{TT} or \textbf{CAL} battery) is \textbf{INVALID}.
\item Unresolved gait-altering pain/mechanics change signals (\textbf{MET-2B-036} = Yes) in the decision window.
\item Unresolved stop-rule events or escalation posture not cleared under governance.
\item Any attempt to salvage gate evidence under unsafe environment conditions rather than reslotting.
\end{itemize}

D) Intended tier posture for Phase 02 exit

\begin{itemize}
\item Primary exit tier posture: \textbf{INTERMEDIATE} (Phase 02 is expected to consolidate an \textbf{INTERMEDIATE} foundation that Phase 03 converts into power/endurance outputs).
\end{itemize}

\subsection*{12. Risk Controls and Stop Rules}
\phantomsection

\begin{itemize}
\item Stage 0.5 and Stage 1.F are controlling authorities for stop posture and escalation. This Phase Card adds phase-specific risk controls without creating a new stop-rule system.
\item \textbf{STR}-specific controls:
\begin{itemize}
\item No true \textbf{1-RM} attempts. \textbf{C17} submax proxy posture is mandatory for any \textbf{STR} proxy evidence.
\item Technique breakdown terminates the exposure; do not “grind through” for numbers.
\item Joint-history conservative constraints remain in force (shoulder/hip/ankle): avoid ballistic end-range stress and unstable loading setups; regress immediately on mechanics change.
\end{itemize}
\item \textbf{RUN}-specific controls:
\begin{itemize}
\item \textbf{RUN} \textbf{TT} (\textbf{C9}) is executed only in protected windows and only under safe conditions.
\item Environment hazards (ice, low visibility, extreme windchill, unsafe heat, poor air quality) trigger reslot or treadmill substitution only when full identity/settings logging preserves comparability; otherwise reslot.
\end{itemize}
\item \textbf{CAL}-specific controls:
\begin{itemize}
\item Artifact evidence is a validity requirement for gate-grade \textbf{CAL} evidence. Missing artifact makes the test \textbf{INVALID} for gate use.
\item Avoid clustering \textbf{CAL} battery too close to \textbf{RUN} \textbf{TT} or \textbf{STR} proxy sessions; respect Stage 3 spacing intent.
\end{itemize}
\item Durability disqualifier controls:
\begin{itemize}
\item \textbf{MET-2B-036} (gait-altering pain flag) = Yes is treated as a gate-disqualifier signal that constrains both scheduling (abort/reslot) and advancement posture.
\item \textbf{MET-2B-037} (foot hotspot/blister) is treated as a primary constraint for \textbf{RUN} scheduling and for any exposure that would worsen skin integrity.
\end{itemize}
\item Recovery drift controls:
\begin{itemize}
\item \textbf{RECOVERY} markers (sleep, fatigue, soreness, session RPE) are used to decide whether a protected week remains test-capable or becomes recovery-first with reslot of major tests.
\end{itemize}
\item Gate-grade behavior controls:
\begin{itemize}
\item No make-up stacking after missed/\textbf{INVALID} tests.
\item No novelty near protected windows (route, tool, footwear, or timer changes) unless forced; forced changes are comparability breaks and bias toward conservative interpretation and/or reslot.
\end{itemize}
\end{itemize}

\clearpage
\begin{landscape}

\subsection*{13. Required Capabilities and Dependencies}
\phantomsection

\begingroup
\scriptsize
\setlength{\tabcolsep}{2pt}
\renewcommand{\arraystretch}{1.12}

\begin{longtblr}{
width=\linewidth,
rowhead=1,
colspec = {
Q[l,wd=0.90in]
X[l,3.05]
X[l,1.55]
Q[l,wd=1.20in]
X[l,3.20]
},
hlines,
vlines,
cells = {hsep=6pt, vsep=6pt, halign=l, valign=t},
cell{1-Z}{1-5} = {fg=black},
row{odd} = {bg=white},
row{even} = {bg=LightGray},
row{1} = {bg=StageBlue!15, fg=black, font=\bfseries, halign=l, valign=t},
}
\Hdr{Dependency \textbf{ID}} &
\Hdr{Required Capability Category and Tier} &
\Hdr{Due-By week-in-phase} &
\Hdr{Owner} &
\Hdr{Fail Action} \\

\textbf{DEP-1C-021} &
7.16 Operational Self-Management Systems — Tier 1 (same-day logging reliability; schedule control; audit-ready entries) &
Before Week 1 (baseline) &
Trainee &
If logging cannot be executed same-day, treat evidence as inadmissible for gate decisions; default \textbf{HOLD} posture for gate use until a 14-day compliance proof exists; reslot any planned gate-grade tests rather than forcing compromised records. \\

\textbf{DEP-1C-006} &
7.18 Testing, Timing, Measurement, and Course-Setup Tools — Tier 1 (reliable timer method; stable route identity; treadmill identity and settings logging capability) &
Before Week 1 (baseline) &
Equipment Lead &
If timer, route, or treadmill identity cannot be logged, \textbf{RUN} \textbf{TT} evidence is \textbf{INVALID} for gate use; reslot \textbf{C9} to next protected window. Do not substitute \textbf{AER} evidence to claim \textbf{RUN} tier attainment. \\

\textbf{DEP-1C-008} &
7.11 Footwear Systems and Foot Care — Tier 1 (stable footwear interface; hotspot management; no-novelty posture near tests) &
Before Week 1 (baseline) and maintained through Week 12 &
Equipment Lead &
If foot hotspots/blisters escalate or gait is altered, downshift exposure and reslot \textbf{RUN} \textbf{TT}; default \textbf{HOLD} for advancement if gait-altering pain flag becomes Yes or if foot integrity is not stable approaching gate windows. \\

\textbf{DEP-1C-012} &
7.01 Strength and Resistance Training Equipment — Tier 1 (safe external loading capability sufficient for \textbf{C17} proxy posture; stable setup method; pull-up interface feasibility) &
By Week 4 Deload+Test &
Equipment Lead &
If safe external loading capability is missing or unstable, \textbf{C17} proxy is not executed and is reslotted to the next protected window; Phase 02 remains training-focused with no gate claims from \textbf{STR} proxies until capability exists. \\

\textbf{DEP-1C-014} &
7.08 Conditioning, Endurance, and Aerobic Training Tools — Tier 1 (reliable low-impact Cardio modality access under Grand Forks conditions) &
Before Week 1 (baseline) &
Coach Lead &
If no safe Cardio modality access exists, sessions cannot remain \textbf{VALID} under global constraints; default posture is \textbf{HOLD} and feasibility intervention until Cardio continuity is achievable without unsafe exposure. \\

\textbf{DEP-1C-019} &
7.14 Environmental Apparel — Tier 1 (minimum capability to execute safe outdoor exposure or shift indoors when unsafe; documentation posture for substitutions) &
Before Week 8 \textbf{RUN} \textbf{TT} window &
Trainee &
If environment hazards prevent safe, comparable \textbf{RUN} testing and indoor substitution cannot preserve validity, reslot the \textbf{RUN} \textbf{TT}; do not force compromised attempts due to schedule pressure. \\

\textbf{DEP-1C-016} &
7.10 Logistics, Maintenance, Storage, and Sustainment Equipment — Tier 1 (binder/log storage, drying/maintenance capability sufficient to prevent repeatable foot breakdown) &
By Week 2 &
Trainee &
If drying/storage is insufficient and foot breakdown trends rise, downshift exposure and treat as durability constraint; default \textbf{HOLD} for advancement if durability flags cannot be stabilized. \\

\textbf{DEP-1C-017} &
7.12 Recovery, Mobility, and Tissue-Care Implements — Tier 1 (minimum mobility/recovery capability to sustain 6 sessions/week under rising strength demand) &
By Week 2 &
Trainee &
If recovery drift escalates (sleep collapse, soreness/fatigue trends), convert protected windows to recovery-first and reslot major tests; repeated drift forces \textbf{HOLD} until stability returns. \\

\end{longtblr}

\endgroup

\subsection*{14. Validity Requirements}
\phantomsection

\begin{itemize}
\item Session evidence:
\begin{itemize}
\item Sessions are tagged \textbf{VALID} / \textbf{MODIFIED} / \textbf{INVALID} per Stage 1.F and must include complete required elements and minimum dataset.
\item Cardio minimum standard met (\textbf{MET-2B-045}) is required for admissible sessions; a session below the minimum standard is \textbf{INVALID} for gate purposes.
\end{itemize}
\item Test evidence:
\begin{itemize}
\item Tests are tagged \textbf{VALID} / \textbf{INVALID} only. \textbf{INVALID} provides no gate credit and is treated as not yet tested for gate decisions until rerun under valid conditions.
\item Gate-grade tests are executed inside protected windows only (Week 1 baseline; Weeks 4/8/12 Deload+Test) per Stage 3.
\end{itemize}
\item \textbf{CAL} validity (\textbf{C4}–\textbf{C7}):
\begin{itemize}
\item Artifact evidence is required for gate-validity. If artifact is missing or non-compliant, the test is \textbf{INVALID} for gate use and must be reslotted.
\end{itemize}
\item \textbf{RUN} validity (\textbf{C9} / \textbf{MET-2B-015}):
\begin{itemize}
\item Route identity must be verified and logged (\textbf{ROUTE} identity) or treadmill identity must be logged (\textbf{TM} identity) with required settings logging.
\item If outdoor conditions are unsafe, treadmill substitution is permitted only under full identity/settings logging; otherwise reslot.
\end{itemize}
\item \textbf{STR} proxy validity (\textbf{C17} / \textbf{MET-2B-027}–031):
\begin{itemize}
\item Implement identity and loading method must be logged.
\item Submax, RPE-capped posture is mandatory; any deviation toward true max posture is treated as invalid evidence for tier claims.
\end{itemize}
\item Disqualifier integration:
\begin{itemize}
\item \textbf{MET-2B-036} (gait-altering pain flag) = Yes is a gate disqualifier signal; major tests are aborted/reslotted rather than forced under compromised mechanics.
\end{itemize}
\end{itemize}

\end{landscape}

\subsection*{15. Change Control Notes}
\phantomsection

Allowed without patch (must be documented and auditable):

\begin{itemize}
\item Swap sequencing of tests within the same protected window week without adding events and without violating cap/spacing intent.
\item Reslot a missed/\textbf{INVALID} test to the next protected window within Phase 02.
\item Convert a protected Deload+Test week into recovery-first or check-in-only due to readiness drift and reslot maximal tests forward.
\item Substitute \textbf{RUN} route to treadmill only when validity/logging preserves comparability; otherwise reslot.
\end{itemize}

Requires patch (Stage 1.F change control posture; revalidation required):

\begin{itemize}
\item Extending Phase 02 duration.
\item Adding a new protected window not present in Stage 3 for Phase 02.
\item Changing the gate test set (adding/removing major gate constructs, changing \textbf{RUN} distance construct, changing \textbf{CAL} battery composition beyond Stage 2).
\item Changing the work:deload ratio or protected window placement pattern.
\item Changing the phase mission/purpose in a way that alters domain priorities or gate posture.
\end{itemize}

Reslot discipline reminder:

\begin{itemize}
\item Missed/invalid tests are corrected by reslot, not by stacking multiple major tests into the same week or violating spacing rules.
\end{itemize}

\subsection*{16. Compliance Checklist}
\phantomsection

\begin{itemize}
\item Pass: Phase \textbf{ID}, name, and duration match the locked phase ledger (Phase 02; 12 weeks).
\item Pass: Operating constraints explicitly state 6 sessions/week, required session elements in correct order, Cardio minimum standard in body text, and 10,000 steps/day global requirement in body text.
\item Pass: Weekly structure table includes Sessions 1–6 and each Cardio Modality/Intent cell contains Modality; Minutes; RPE plus Talk Test; Continuity statement.
\item Pass: Deload/test placement aligns to Phase 02 protected windows (Week 1 baseline; Weeks 4/8/12 Deload+Test).
\item Pass: Test Battery table references Stage 2 instruments only (\textbf{C4}–\textbf{C7}, \textbf{C9}, and \textbf{C17} only if finalized and deployable, with \textbf{MET-2B} IDs) and does not invent tests.
\item Pass: Exit standards enforce Stage 1.F admissibility rules, including 5-of-6 admissible sessions per week on average, \textbf{MODIFIED} cap posture, and the binary disqualifier that any \textbf{INVALID} session in the decision window disqualifies \textbf{ADVANCE}.
\item Pass: Dependencies table uses Stage 7 category IDs 7.01–7.20 with tier posture, due-by timing, owners, and fail actions that default to \textbf{HOLD}/reslot rather than unsafe forcing; 7.02 is not required.
\item Pass: Governing Status must remain Draft — Non-Deployable until \textbf{C17} is fully finalized and non-placeholder.
\end{itemize}

\end{document}